\section{The finite Whittaker test and Plancherel identity}
\label{sec:plancherel}

We complete the deformation argument of Section~\ref{sec:deformation}.
Its remaining input is a uniform upper weight bound for the finite
approximations to the cone of \eqref{eq:P6}.  A finite Whittaker test makes
that cone vanish at exponent zero and detects its vanishing at irrational
exponents.  The weight bound will then allow Lemma~\ref{deform:colimit} to
transfer the vanishing.  Throughout, $Y=Y_H$, $b=\dim Y$, and
$D_s=D(L_H^s)$, as in that section.  All the tests below act in the output
bundle variable; the spectator variable is retained.

\subsection{A test at one point}
\label{planch:finite-test-construction}

Fix $x_0\in X$ and a parameter $z$ at $x_0$.  Set
\[
 K=L^+H,\qquad K_1=\ker(K\longrightarrow H),\qquad
 I^-=\operatorname{ev}^{-1}(B^-),
\]
and define
\[
 K^d=\operatorname{Ad}_{z^{-\rho}}K,\qquad
 I^d=\operatorname{Ad}_{z^{-\rho}}I^-,\qquad
 K^d_1=\operatorname{Ad}_{z^{-\rho}}K_1.
\]
These conjugations take place in the adjoint group.  They do not require
$\rho$ to belong to the cocharacter lattice of $H$.
Write $Y(J)$ for the stack with level $J$ at $x_0$, so that $Y(K)=Y$.
One may construct these stacks by formal frames and disk gluing, or by the
corresponding smooth parahoric group schemes.  The level changes used here
are finite-type morphisms: after passage to a sufficiently deep congruence
subgroup, they are the usual finite-dimensional homogeneous-space
fibrations of jet-frame torsors.  In particular, with
\[
 V=Y(K^d_1),
\]
the map $V\to Y(K^d)$ is an $H$-torsor and $V/B^-=Y(I^d)$.

We define a functor $\mathcal T_s$ in three steps.  First pull by $!$ from
$Y$ to $Y(I^-)$.  Next apply the star intertwiner
\[
 p_{2,*}p_1^!,\qquad
 Y(I^-)\xleftarrow{\ p_1\ }Y(I^-\cap I^d)
             \xrightarrow{\ p_2\ }Y(I^d).
\]
Finally apply finite $(N,\chi_0)$ star averaging, where $\chi_0$ is a
nondegenerate residual character.  Either sign of $\chi_0$ is allowed.
The last operation is the kernel transform for the open incidence
\begin{equation}
 \{([vn],v):v\in V,\ n\in N\}
 \ \subset\ V/B^-\mathbin{\times}_{Y(K^d)}V:
 \qquad ([vn],v)\longmapsto \chi_0(n).
 \label{eq:P7}
\end{equation}
It uses $!$-pull from the first projection, tensor with the exponential of
the displayed function, and star image to $V$.  Fixed normalization shifts
do not affect any vanishing or conservativity assertion in this
subsection.

Here are the twisting lines in this construction.  On $Y(K^d)$ use the
determinant line defined by the lattice
$\operatorname{Ad}_{z^{-\rho}}\mathfrak h_{\mathcal O}$, where
$\mathcal O=\mathbb C[[z]]$.  On $Y(I^d)$ correct its pullback by the
inverse relative determinant of this lattice and
$\mathfrak h_{\mathcal O}$.  On frames the relative determinant is a fixed
vector line whose torus character is
\begin{equation}
 \nu=\sum_{\alpha>0}2\operatorname{ht}(\alpha)\alpha
       =\operatorname{Kil}(\rho,-).
 \label{planch:det-character}
\end{equation}
The formula follows by counting the shifted root-lattice lines.  Extend
this character to the Iwahori through its torus quotient and form the
associated line.  Its inverse is the required correction.  After taking
exponent $s$, the resulting line agrees with the input twisting line on
the intersection correspondence.  Characters of that intersection are
determined on its torus, so this check proves the agreement globally.
On $V$ the frame trivializes the correction line, with its specified
character law; use the same trivialization at $vn$ in \eqref{eq:P7}.
Thus $\mathcal T_s$ is a well-defined functor from $D_s$ to twisted
D-modules on $V$.

\begin{proposition}[Finite Whittaker test]
\label{planch:finite-test}
The functor $\mathcal T_0$ kills the anti-tempered subcategory of $D(Y)$.
If $s\notin\mathbb Q$, the functor $\mathcal T_s$ is conservative on all
of $D_s$.  Both assertions remain valid after tensoring with a spectator
D-module category.
\end{proposition}

\begin{proof}
For the assertion at zero we use the local Whittaker description of
anti-tempering.  In the notation of
\cite[Lemmas~4.2.0.1 and~4.3.0.1, Equation~(4.3.1)]{FaergemanRaskin},
the kernel of nondegenerate Whittaker $!$-averaging on positive-loop
invariants is exactly the anti-tempered subcategory.  That averaging
factors through the translated first adolescent Whittaker level, whose
inclusion in the full Whittaker category is fully faithful.  Its
finite-level formula is forgetful passage to the opposite Iwahori,
star intertwinement, and residual finite Whittaker averaging.  The last
$!$-averaging agrees with star averaging on these inputs up to a fixed
shift; see also \cite[Section~3.1, Lemma~3.1.4]{BogdanovaWhit}.
The resulting finite-level diagram on bundle stacks is exactly the
construction above.  The use of adjoint conjugation for a group on which
$\rho$ is not a cocharacter is allowed in this theorem; see
\cite[Remark~4.4.0.1]{FaergemanRaskin}.

For completeness, the categorical mechanism explains why this assertion
applies to all D-modules on the level tower.  An object of positive-loop
invariants determines an equivariant functor out of the Grassmannian
module category.  The Whittaker Grassmannian category is a tempered
spherical module.  Consequently the functor associated to an
anti-tempered object becomes zero after taking Whittaker coinvariants:
the tempered projector acts as the identity on that module, whereas
precomposition with the projector kills the given functor.  Whittaker
invariants and coinvariants agree in the character conventions under
consideration.  The finite test lands in the first adolescent level,
where the fully faithful inclusion detects zero.  Its operations commute
with the equivariant functor: this is clear for restriction and finite
averaging, and for the intertwiner follows from the finite unipotent
root-piece quotients of the two compact subgroups by their intersection.
This also proves the assertion with spectators, without a holonomicity
restriction.

We next prove conservativity for irrational $s$.  The first $!$-pull is
conservative because the level map is smooth and surjective.  To treat
the intertwiner, choose a minimal gallery from the alcove of $I^-$ to
that of $I^d$.  The open Bruhat convolution for the associated reduced
word identifies its product correspondence with the intersection-level
correspondence.  The star intertwiner therefore factors into the
intertwiners for adjacent alcoves.  Twists factor as well: at each
Iwahori use its own Lie-lattice determinant and the inverse relative
determinant with respect to $\mathfrak h_{\mathcal O}$.  Transitivity of
relative determinant gives the comparisons at consecutive intersections
and recovers the prescribed endpoint lines.

For adjacent alcoves, work smoothly locally over the wall-level stack.
The correspondence is the open of distinct flags in a product of two
projective lines for the rank-one reductive quotient.  The wall-lattice
determinant comes from the base; the remaining twists are the torus lines
on the two flag varieties.  Compose the star open kernel with the
reverse shriek open kernel.  To calculate this composition, fix the
first and last flags.  If they differ, the middle flag ranges over a
twice-punctured projective line, with star extension at one puncture and
shriek extension at the other.  The regular-singular rank-one
connection occurring in the tensor product has the same cohomology as
its restriction to the link of the shriek puncture; the relative
cohomology specified by shriek extension at that puncture is
therefore zero.  In coordinates with punctures $a\ne c$, this is the
fibre of the restriction map
\[
 R\Gamma(\mathbb P^1\setminus\{a,c\},\mathcal L)
       \longrightarrow R\Gamma(\operatorname{link}_c,\mathcal L),
\]
which is an isomorphism for every Kummer system $\mathcal L$, including
the trivial one.  The $!$-tensor calculation has this description by
Verdier duality and the projection formula.  If the flags coincide,
the ratio twists cancel.  The
projection formula for star image and $!$-tensor reduces the calculation
to the constant kernel on an affine line, whose integral is a nonzero
constant complex up to shift.  The calculation is smooth along the
diagonal.  The composite is thus a nonzero constant multiple, up to
shift, of the identity kernel.  Each adjacent intertwiner is
conservative.  The alcove argument uses affine Weyl representatives,
which exist for the actual root datum; it does not assert the existence
of $z^{-\rho}$ in $LH$.

It remains to treat finite Whittaker averaging.  Smoothly locally over
$Y(K^d)$, its source is the flag variety $H/B^-$ and its output is the
frame space $H$.  The flag twist has torus character $-s\nu$.
Compose its star open kernel with the reverse star open kernel, using
opposite pairing twists and the negative exponential.  By the
projection formula this is computed on the two open incidences.
The composite is equivariant for the left $H$-action on the pair of
flags.  On a Bruhat stratum indexed by $w\ne1$, the stabilizer torus
carries the Kummer character
\[
 s(w\nu-\nu)
\]
up to the harmless overall sign convention.  The weight $\nu$ is regular
on the root system.  Thus $w\nu-\nu$ is nonzero, and irrationality of
$s$ makes the character nonintegral on some cocharacter of the
stabilizer torus.  The Kummer inertia test forces the equivariant
category on that stratum to vanish.

For the diagonal calculation we use $!$-restriction, rather than
ordinary star pull.  If $j_1,j_2$ are the two open incidences in the
space of a pair of flags and an intermediate frame, the $!$-tensor
projection formula and $!$-base change identify the tensor of their
star extensions with star image from their open intersection.
For the projection $q$ forgetting the intermediate frame and the
diagonal inclusion $i_\Delta$, schematic $!$-base change gives
$i_\Delta^!q_*=q_{\Delta,*}(i'_\Delta)^!$, with $i'_\Delta$ the pulled-back
diagonal.  Thus no interchange of nonproper star image with ordinary
star restriction is being used.
On the diagonal the two $n$-coordinates agree.  Both the exponential
factors and the twisting factors cancel.  After fixing a flag and a
representative of it, the remaining integral is the de Rham integral of
the dualizing object on $N\times B^-$.  Its cohomology is nonzero and
varies constantly along the diagonal.  The composite therefore acts as
the identity tensored with a nonzero complex of vector spaces, and is
conservative.  These are ordinary D-module kernel calculations on fixed
finite-dimensional correspondences.  They apply to all D-modules and
commute with spectator factors.  This proves the proposition.
\end{proof}

\subsection{Preparing both character parameters}
\label{planch:two-axes}

Let $\{\widehat C_\beta\}$ be the finite augmented cone diagrams
constructed at the end of Section~\ref{deform:finite-stages}.  They
already carry the pre-axis $u$.  Write $\widehat{\mathcal T}$ for the
pre-axis version of the finite Whittaker test: adjoin
$\delta_{u'=0}$, perform its pullback and intertwiner, and replace the
exponential in \eqref{eq:P7} by translation of the second axis $u'$ by
$\chi_0(n)$.  Fourier specialization at its character parameter $1$
recovers $\mathcal T_s$.

The new finite diagram has the dilation property in both axes.  For
$u'$ use the right torus action on frames: it conjugates $n$, scales
$\chi_0$ and $u'$ by the same square, and leaves the first flag $[vn]$
fixed.  The correction line transforms by its character law.  The
matched determinant lines define maps between frame torsors, and
$\widehat{\mathcal T}$ transports the single coefficient input along
these maps.  At logarithm length $a$ this is one $\mathcal L_a$ input;
no tensor product over $\mathbb Q$ of separate logarithm variations is
taken.  A change of trivialization merely rescales the one frame-ratio
coordinate.  Thus these operations introduce no new monodromy input.

Define the tested cones
\begin{equation}\label{planch:tested-cones}
 \widehat K_\beta
   =C_!^u C_!^{u'}\widehat{\mathcal T}(\widehat C_\beta).
\end{equation}
On fixed output and spectator charts, including their frame torsors,
these are finite geometric calculations linear in the original Kummer
input, with the compatible logarithm lifts of
Section~\ref{sec:deformation}.  To apply Lemma~\ref{deform:colimit} we
must prove two properties of this family: its cohomology colimit is zero
at $q=1$, and its weights in each stalk degree have an upper bound
independent of $\beta$.  The next lemma proves the first property; the
rest of the section establishes the bound.

\begin{lemma}
\label{planch:zero-colimit}
For every stalk $\xi$ on a test chart and every degree $j$,
\[
 \operatorname*{colim}_\beta
       H^j\bigl((\widehat K_\beta)_{\xi,q=1}\bigr)=0.
\]
\end{lemma}

\begin{proof}
With both Fourier parameters nonzero, the exponent-zero Poincar\'e
counit has anti-tempered cone by the enhanced classical Whittaker
theorem recalled in Section~\ref{sec:deformation}.  Proposition~
\ref{planch:finite-test} kills this cone.  Continuity gives the same
statement with spectators.  Dilation spreads vanishing at unit
parameters to the product of the two punctured Fourier axes.

This argument also commutes with the colimit used here.  On smooth
charts the calculations lie in Ind-holonomic complexes.  The dilation
identifications hold stage by stage, so the colimit vanishes after
pullback to the dilation covers and hence before that pullback.
Shriek extension from the punctured axes is computed in the same
Ind-holonomic category and preserves this zero colimit.  The Fourier
description of $C_!^u C_!^{u'}$ now gives the assertion for the
tested cones $\widehat K_\beta$.
\end{proof}

We will need one boundedness property of the fixed finite test.  At
$q=1$, its pre-axis version increases upper weights and upper ordinary
cohomological degrees by fixed constants on fixed test charts.  To see
the weight assertion, factor the intertwiner into adjacent-alcove
steps and compactify each open by its proper wall-flag correspondence.
Smoothly locally over the input flag, the star extension is external
product with a fixed fibre kernel.  Its weight cost is therefore fixed;
the remaining push is proper.  The same reasoning applies to
\eqref{eq:P7}: compactify the incidence inside
$V/B^-\times_{Y(K^d)}V$, retaining $u'$.  To see the product over the
input, choose a local lift $g(f)$ of its flag $f$.  The incidence
$[vn]=f$ writes $v=g(f)b n^{-1}$, with $b\in B^-$.  In the coordinate
$g(f)^{-1}v$, its open fibre is the fixed cell $B^-N$ and its character
graph $u'=\chi_0(n)$ is independent of $f$.  An arbitrary input is
therefore external-tensored with this fixed star-extended graph kernel.
The remaining projection, from the compactification to $V$ with the
axis retained, is proper.  Ordinary cohomological dimension gives the
corresponding fixed degree bounds.  Above a fixed output chart these
level and flag correspondences require only fixed quasi-compact input
charts.  Finally each pre-axis projection uses a one-dimensional
projection with fixed shifts.  None of these costs depends on the
number of marks, the pole bound, the Satake labels, or the skeletal
degree.

The uniform estimate has three parts.  The two opposite Hecke kernels
have cancelling weights, removing dependence on the Satake labels.
A spectral calculation then gives an upper cohomological degree
independent of those labels and the number of marked points.
Configuration integration lowers this degree by the collision
codimension, and the simplicial shifts lower it further.  Only finitely
many collision codimensions and simplicial degrees can therefore
contribute to any fixed stalk degree; their weight bounds will suffice.

\subsection{Satake decomposition of the finite terms}
\label{planch:satake-decomposition}

All weights from now until the final application are computed at
$q=1$.  On frame-torsor charts the monodromy input is then constant in
mixed Hodge modules.  We use the mixed-complex convention of
Section~\ref{sec:deformation}, and work on fixed final test charts.
Smooth-chart dimension shifts are fixed constants there.

Take a term $\pi_!\widehat M_{I,m}$ of \eqref{eq:P4}, retaining its
pre-axis.  Stratify $X^I$ by coincidence patterns.  The filtration by
$!$-restrictions expresses it by star extensions from these strata.
Let $S_\alpha$ have $n$ distinct points.  Its closure is a diagonal copy
of $X^n$, denoted $\overline S_\alpha$, and
$j_\alpha:S_\alpha\hookrightarrow\overline S_\alpha$ is the open
configuration inclusion.  We first retain the marked parameters and
use the entire stack $Y$ as spectator, with
$F_Y=\Delta_!k_Y$ in output--spectator order.  All comparisons are
relative and will subsequently
be pulled to a smooth spectator chart $W$.  Comparisons made modulo
$u$-constants are used only after applying the pre-axis projections.

For a tuple $\lambda=(\lambda_1,\ldots,\lambda_n)$ of dominant
coweights within the pole bound, let $\mathsf S_\lambda$ be the relative
Satake operator.  Our convention on bounded holonomic diagrams is
\[
 \mathsf S_\lambda(D)
   =p_{2,!}\bigl(p_1^*D\otimes^* K_\lambda\bigr),
\]
where $p_1$ is the input projection and the Hecke projections are proper
on the support.  The labels are chosen so that application to the
vacuum gives the corresponding Whittaker standard objects.  Extend
the kernels over $\overline S_\alpha=X^n$ by fusion.  On the bundle
Hecke correspondence set
\[
 d_\lambda=\sum_i\langle2\rho_{\rm rt},\lambda_i\rangle,
 \qquad
 K_\lambda=\operatorname{IC}_\lambda[-(b+n)].
\]
There are no half Tate twists.  Thus $K_\lambda$ has complex weight
$d_\lambda$, and its relative dual over the smooth base of dimension
$b+n$ is
\begin{equation}
 K_\lambda^r=K_\lambda(d_\lambda),\qquad
 \operatorname{wt}(K_\lambda^r)=-d_\lambda.
 \label{planch:dual-satake}
\end{equation}
The right adjoint $\mathsf S_\lambda^r$ uses the exchanged paths and
this relative dual.  The Tate twist disappears after forgetting the
mixed structure, but its weight in \eqref{planch:dual-satake} is
essential below.

We recall the exchange properties being used.  Fusion Satake kernels
are universally locally acyclic (ULA) relative to either Hecke
projection.  The relative duality formula identifies
$K\otimes^*p_i^*(-)$ with relative Hom from the dual kernel into
$p_i^!(-)$; kernels and their relative duals commute with base change.
Together with properness, this gives the displayed adjoints, exchange
with $!$-pull, and exchange with star image under parameter changes.
These are comparisons in mixed complexes.  They may be checked after
forgetting to constructible complexes, because the comparison maps,
duals, and adjunctions are themselves mixed.  At distinct fixed points
the Hecke correspondences are locally products over both projections,
using finite jet frames.  Thus the same formulas hold for the
holonomic exponential expressions after Fourier specialization.
We use the usual fusion and Casselman--Shalika theorems with these
relative normalizations; see \cite{BravermanGaitsgory,GLCII}.

The pole-model Hecke correspondence is the base change of the bundle
Hecke correspondence along either projection: modifying the extension
at the marked points transports the Pl\"ucker data.  Its kernels are
the star pullbacks of $K_\lambda$.  All calculations in a fixed bound
use finite bounds.  In particular, a modification bounded by
$\lambda$ takes the vacuum pole bound to the bound allowing exactly
those dominant pole orders on the highest-line maps.  This is the
Schubert-bound estimate on representation lattices; at collisions
the pole orders add.  It also permits the adjoint calculation back
to the vacuum using just the indicated forward bound.

Over $S_\alpha$, denote the vacuum operations of \eqref{eq:P3} by
$\widehat T_0$ and $\widehat H_0$.  Their vacuum stratum consists of
true $N_0$-bundles, with no poles or defects.  Put
\[
 \widehat Q_\lambda=\mathsf S_\lambda\widehat T_0,
 \qquad
 \widehat H^\lambda=\widehat H_0\mathsf S_\lambda^r.
\]
The following arrows are the evaluation maps of these adjunctions and
the map induced by the Whittaker projection counit:
\begin{equation}
 \widehat M\longleftarrow
 \bigoplus_\lambda\widehat Q_\lambda
                    \widehat H^\lambda(\widehat M)
 \longrightarrow
 \bigoplus_\lambda\widehat Q_\lambda
                    \widehat H^\lambda(\widehat F).
 \label{eq:P8}
\end{equation}
Here $\widehat M$ and $\widehat F$ are $!$-restricted to $S_\alpha$,
and the finite sum ranges over its labels in the pole bound.  Both
arrows are isomorphisms modulo $u$-constants.  Indeed Fourier commutes
with the Hecke operations, whose kernels are external to the axis.
The diagram has dilation because Hecke transport commutes with the
Cartan action on Pl\"ucker data.  At a nonzero character parameter,
Casselman--Shalika identifies the vacuum translates with the standard
objects of the Whittaker summands, while $\widehat M$ is the Whittaker
projection.  This gives the assertion.  The verification may be made
on $!$-fibres over $S_\alpha$, also with spectators: properness and
the local product description give the Hecke base changes, the
Whittaker decomposition gives the projection base changes, and the
vacuum bound is constant over the parameter base.  Notice that the
restricted $\widehat F$ uses $\pi^!$ over $S_\alpha$, so includes the
base dualizing factor.

Define two objects without marked parameters:
\begin{equation}
 U_0=\pi_!\widehat T_0(\delta_{u=0}),\qquad
 U'_0=\widehat H_0
       \bigl(\pi^!(F_Y\boxtimes\delta_{u=0})\bigr).
 \label{eq:P9}
\end{equation}
They live respectively on the output and spectator copies of $Y$,
each with a pre-axis.  Equivalently one may use a fixed marked point
with pole bound zero in this definition.  Write $*_{!,u}$ for external
product followed by $!$-image under addition of the two pre-axes.

\begin{lemma}[Separation of the two bundle variables]
\label{planch:separation}
After \eqref{eq:P8}, the term indexed by $\lambda$, star extended to
$\overline S_\alpha$ and pushed by $\pi_!$ with marked parameters
retained, is, modulo $u$-constants,
\begin{equation}
 \mathsf S_\lambda^{(1)}\mathsf S_\lambda^{\prime(2)}
 \left((U_0*_{!,u}U'_0)\boxtimes
                    (j_\alpha)_*\omega_{S_\alpha}\right).
 \label{eq:P10}
\end{equation}
Both Hecke operations are over $\overline S_\alpha$.  The second is
the transpose-path version of $\mathsf S_\lambda^r$: its path is
reversed back to the forward direction, but its kernel is
$K_\lambda^r$, rather than $K_\lambda$.  Formula~\eqref{eq:P10}
includes the base dualizing factors and has no additional shifts
depending on $I$, $\alpha$, or $\lambda$.
\end{lemma}

\begin{proof}
First work over $S_\alpha$.  In
$\widehat H^\lambda(\widehat F)$, move
$\mathsf S_\lambda^r$ from the pole model to the bundle stack.
The Cartesian Hecke squares, properness, and the ULA exchange formulas
justify this move through $\pi^!$.  Only the vacuum restriction for
$\widehat H_0$ is needed, so the forward Schubert bound just described
suffices.  Acting on the first slot of the $!$-diagonal kernel is the
same as acting on the second slot by the transpose-path operator
$\mathsf S'_\lambda$.  In both expressions, star pull followed by
the diagonal $!$-image produces the same $!$-image along the Hecke
path, with the same relative dual kernel.  Now commute
$\widehat H_0$ with that spectator operation.  Its $!$-restrictions,
vacuum translation, and star image satisfy the same exchange
formulas.  This gives
$\mathsf S'_\lambda(U'_0\boxtimes\omega_{S_\alpha})$.

Next apply $\widehat T_0$ and the forward pole-model Hecke operator.
The vacuum operation commutes with the spectator Hecke calculation
by star-pull and shriek calculus.  Before either Hecke operation we
therefore have an external object with parameter factor
$\omega_{S_\alpha}$.  The ULA fusion kernels exchange with
$(j_\alpha)_*$, including the closed embeddings of pole bounds.
Consequently extension to $\overline S_\alpha$ changes only that
external factor, to $(j_\alpha)_*\omega_{S_\alpha}$.

Finally move $\pi_!$, with parameters retained, through the
pole-model Hecke operation.  The Cartesian square and the exchange
between star pull and $!$-image turn it into bundle-stack Hecke.
Before this operation, the vacuum translation formula is exactly
\[
 \widehat T_0(E)
   =\widehat T_0(\delta_0)*_{!,u}E.
\]
It gives the addition convolution in \eqref{eq:P10}.  Every step
is relative to the spectator and hence commutes with smooth pull to
$W$.  The closed diagonal embedding
$\overline S_\alpha\hookrightarrow X^I$ may be included afterwards
and introduces no further shifts.
\end{proof}

\subsection{Uniform bounds before configuration integration}
\label{planch:uniform-bounds}

Replace the two vacuum inputs by
\[
 \bar U_0=C_!^uU_0,\qquad \bar U'_0=C_!^uU'_0.
\]
This does not change \eqref{eq:P10} modulo $u$-constants: addition
convolution with a $u$-constant is again $u$-constant, by a change
of coordinates and the $!$-projection formula.  The benefit is that
both bar objects have a global upper weight bound, even though $Y$
is not quasi-compact.

\begin{lemma}
\label{planch:vacuum-weight}
The objects $\bar U_0$ and $\bar U'_0$ have fixed upper complex-weight
bounds on their entire bundle stacks.  In \eqref{eq:P10}, omit
$(j_\alpha)_*\omega_{S_\alpha}$ and instead use the unshifted constant
complex $k_{\overline S_\alpha}$, of weight zero and ordinary degree
zero.  With the bar inputs, the resulting
object has an upper complex-weight bound independent of
$I,\alpha,\lambda$ and the pole bound.
\end{lemma}

\begin{proof}
For $U_0$, the vacuum $!$-image factors through shriek extension from
a fixed quasi-compact open of $Y$.  Its source is the finite-type
stack of $N_0$-bundles with fixed Cartan bundle, and its image is
contained in such an open.  The mixed object on that open is bounded;
shriek extension preserves upper weights.  Applying $C_!^u$ has a
fixed weight cost.

For $\bar U'_0$, use Fourier only to identify support, not to assert
preservation of Hodge weights.  At nonzero Fourier parameter,
$\mathsf F U'_0$ is the $!$-Poincar\'e object of the opposite vacuum,
up to a fixed shift and line.  This is Lin's vacuum comparison
\cite[Theorem~9.0.2 and Section~8.1]{Lin}, applied to the coefficient
of the $!$-diagonal kernel as in \eqref{eq:W2} and \eqref{eq:P5}.
More explicitly that coefficient is
\[
 q_{0,*}\bigl(\operatorname{Exp}(-\psi_0)\otimes^!
                  f_0^!\pi^!F_Y\bigr)[-2\delta_N].
\]
Dilation changes the character parameter while leaving the bundle
variable fixed.  Hence over the punctured Fourier axis this object
extends by shriek from one fixed quasi-compact bundle open.
The Fourier transform of $\bar U'_0$ is its shriek extension across
the zero parameter.  Inverse Fourier, which operates only in the
axis variable, preserves the same bundle-open support property.
The shriek-extension comparison exists in mixed complexes and is an
isomorphism after forgetting the mixed structure.  It is therefore
an isomorphism in mixed complexes.  On the fixed quasi-compact open,
$\bar U'_0$ has an upper weight bound, which extends to all of $Y$.

Apply the two Hecke operations to these bar inputs.  Star pull,
star tensor, and proper image give upper weight bounds, and the
kernel contributions $d_\lambda$ and $-d_\lambda$ cancel by
\eqref{planch:dual-satake}.  Addition convolution uses $!$-image and
has a fixed cost.  Thus the upper bound is independent of every
index.  Restoring the parameter factor in \eqref{eq:P10} simply
star-tensors this object with its pullback, by projection formula.
\end{proof}

Restoring $(j_\alpha)_*\omega_{S_\alpha}$ will introduce weight bounds
that depend on collision codimension, which is unbounded as the number
of points grows.  We now prove a uniform degree bound that will limit
the codimensions contributing to a fixed cohomology degree.  This is
the one place where the full classical geometric Langlands theorem
enters the estimate.

\begin{lemma}
\label{planch:satake-degree}
Fix smooth output charts on the two copies of $Y$.  With the bar inputs
and the parameter factor $k_{\overline S_\alpha}$ used in
Lemma~\ref{planch:vacuum-weight}, the star restriction of
\eqref{eq:P10} to every fibre of $\overline S_\alpha$ has a uniform
upper ordinary cohomological degree.  The bound includes all tuples
of points, all further collisions, and all Satake labels.
\end{lemma}

\begin{proof}
On a parameter fibre, fusion identifies Hecke with fixed-point
Satake operators for tensor products of the given representations.
We first bound the Fourier transform away from zero.  For
$\bar U'_0$, the value at character parameter $1$ on a spectator
chart $W$ is a kernel-valued Hom of the form
\begin{equation}
 \operatorname{Hom}_{D(Y)}
       \bigl(\mathsf S P\operatorname{vac}_!,F_W\bigr),
 \qquad P\operatorname{vac}_!=\pi_!T_0(k).
 \label{eq:P11}
\end{equation}
Here $T_0$ is the exponential-specialized vacuum of the corresponding
sign, $\mathsf S$ is the usual fixed-point Satake operator, and Hom
means application in the first kernel variable, with values in
$D(W)$.  Formula~\eqref{eq:P11} follows from \eqref{eq:P9} by the
transpose calculation in Lemma~\ref{planch:separation}, the smooth
vacuum adjunction with its specified shift, and $!$-image adjunction.
These adjunctions tensor with the spectator.  Their fixed-point
kernel formulas hold for every representation tensor produced by
fusion, with no additional shifts depending on its labels.  Tate
factors can now be forgotten, since we are proving a degree bound.

The object $F_W$ is compact in $D(Y)\otimes D(W)$.  Indeed it is the
$!$-image of a holonomic object on the finite-type affine scheme $W$
along its schematic graph map.  The partial $!$-image has genuine
adjunction against arbitrary D-modules, and the corresponding
$!$-pullback is continuous.  Compactness follows.  The tensor-product
theorem for the ordinary D-module category of $Y$ then expresses
$F_W$ as a retract of a finite construction by extensions and shifts
from exterior products of compact objects.

Let $\mathcal L=\operatorname{LS}_{\check H}$ be the derived stack
of de Rham local systems.  Under classical geometric Langlands,
those first-factor compact objects become coherent compacts in
$\operatorname{IndCoh}_{\operatorname{Nilp}}(\mathcal L)$, and
\[
 \mathsf S P\operatorname{vac}_!
 \quad\longleftrightarrow\quad
 \Xi(\mathcal V),
\]
where $\mathcal V$ is the evaluation representation vector bundle,
or the tensor product of such bundles at the specified points.
The global shift and line are fixed.  This normalization is the
vacuum normalization of \cite[Sections~1.3--1.4]{GLCI}, with Hecke
compatibility from its Theorem~1.2.4 and Proposition~1.7.2, and the full equivalence from
\cite{GLCV}.  More explicitly,
$\mathbb L^L_{H,\operatorname{temp}}(\mathcal O)
=P\operatorname{vac}_!$ and
$\Psi\mathbb L_H=\mathbb L_{H,\operatorname{coarse}}$ by
\cite[Corollary~1.6.5]{GLCI}.  Uniqueness of left adjoints after the
equivalence gives
$\mathbb L_H\mathbb L^L_{H,\operatorname{temp}}=\Xi$.
It follows that the vacuum maps to $\Xi(\mathcal O)$ and its
Satake translates to $\Xi(\mathcal V)$.
The chosen theta characteristic transports the ordinary-D Hecke
conventions; any resulting constant label lines do not change
cohomological degrees.

For a fixed coherent compact $C$, adjunction computes the mapping
complex as
\begin{equation}
 \operatorname{Hom}\bigl(\Xi(\mathcal V),C\bigr)
   =R\Gamma\bigl(\mathcal L,
                    \mathcal V^\vee\otimes^*\Psi(C)\bigr).
 \label{planch:spectral-bound}
\end{equation}
The right adjoint $\Psi$ takes $C$ to a bounded coherent
quasi-coherent complex.  Tensoring with a vector bundle is t-exact,
irrespective of its rank, highest weight, or evaluation points.
The fixed derived QCA stack $\mathcal L$ has finite cohomological
dimension for quasi-coherent global sections in characteristic
zero; we use the QCA finiteness theorem
\cite{DrinfeldGaitsgoryFiniteness}.  Thus
\eqref{planch:spectral-bound} has an upper bound depending on $C$
and $\mathcal L$ only.  Finite extensions, shifts, retracts, and
the fixed compact $D(W)$ factors yield a uniform perverse-D upper
bound in \eqref{eq:P11}.

For $\bar U_0$, Lin's vacuum comparison identifies its Fourier
transform on the puncture with the corresponding transform for
$\bar U'_0$ of the opposite sign, up to a fixed normalization.
Applying fixed-point Satake therefore gives the same bound.  The two
bundle charts may be chosen independently.

The fixed-fibre Satake operations commute with shriek extension
across the Fourier origin, by their kernel formulas.  On the
puncture, dilation makes the preceding bound uniform in the
parameter, with only a fixed dimension shift after the dilation
cover.  The principal affine-open extension $j_{t,!}$ is t-exact
on holonomic complexes.  Fourier is t-exact with its fixed
normalization shift.  The same perverse bounds thus hold on the
whole pre-axis.  The chart dimensions are fixed, so they also give
uniform upper ordinary degrees.  Finally the addition convolution
$*_{!,u}$ has fixed relative dimension and commutes with the
fixed-fibre Hecke calculations.  It changes this upper bound by a
fixed constant.  This proves the lemma.
\end{proof}

\subsection{Configuration strata and the uniform weight estimate}
\label{planch:configuration-estimate}

We now restore $(j_\alpha)_*\omega_{S_\alpha}$ and integrate the
marked parameters.  Stratify $\overline S_\alpha=X^n$ again by
coincidence patterns.  If a stratum $S'$ has $n'$ distinct points,
write
\[
 c'=n-n'\geq0.
\]
In this second filtration we use star restrictions to strata and
shriek extensions for integration.

The star restriction of $(j_\alpha)_*\omega_{S_\alpha}$ to $S'$
has ordinary upper degree
\begin{equation}
 c'-2n.
 \label{planch:link-degree}
\end{equation}
Indeed local coordinates on the curve identify a normal slice with
a product of central braid arrangements of total rank $c'$.  The
local links have the cohomology of their arrangement complements,
whose cohomology vanishes above $c'$.  The dualizing normalization
$[2n](n)$ gives \eqref{planch:link-degree}.
The same restriction has a complex-weight upper bound depending
only on $c'$.  This follows from the mixed arrangement description,
which is \'etale local in curve coordinates.  For fixed $c'$ there
are only finitely many nontrivial collision types: every block of
size $r>1$ consumes $r-1$ units of codimension, and singleton blocks
contribute only smooth factors.  Finally $[2n](n)$ is neutral for
the complex-weight convention.  Thus neither the number of singleton
blocks nor $n$ changes this weight bound.

Let $E$ be the fixed upper degree from
Lemma~\ref{planch:satake-degree}, with the fixed costs of convolution
included.  Integrating the piece over $S'$ by $!$ adds at most
$2n'$ to ordinary cohomological degree, while preserving upper
complex weights.  Its contribution on the test charts therefore
has upper degree
\begin{equation}
 E+c'-2n+2n'=E-c',
 \label{planch:codimension-decay}
\end{equation}
and an upper complex-weight bound depending only on $c'$.
The closed embedding of $\overline S_\alpha$ back into $X^I$ adds
no further cost.

The comparison with the original terms is made after the pre-axis
projections.  Every arrow used in
Lemmas~\ref{planch:separation}--\ref{planch:satake-degree} is
defined in mixed complexes.  Where its being an isomorphism was
checked by Fourier and projection, that check was on underlying
field complexes and therefore also detects the mixed isomorphism.
The fixed finite test and both pre-axis projections add only the
fixed costs established in Section~\ref{planch:two-axes}.
The augmentation-target term has fixed bounds as well.

\begin{proposition}[Uniform weight estimate]
\label{planch:weight-estimate}
Fix a test chart, a stalk $\xi$, and a degree $j$.  The mixed Hodge
structures $H^j((\widehat K_\beta)_{\xi,q=1})$ have a common upper
weight bound, independent of $\beta$ and hence of all pole,
configuration, label, and skeletal indices.
\end{proposition}

\begin{proof}
For skeletal degree $l$, the shift $[l]$ in the finite approximation
lowers \eqref{planch:codimension-decay} to
\[
 E'-c'-l,
\]
where $E'$ includes all fixed shifts.  The complex-weight bound
depends only on $(c',l)$.  A term can contribute to degree $j$
only if
\[
 c'+l\leq E'-j.
\]
Since $c',l$ are nonnegative integers, only finitely many such
pairs occur.  Take the maximum of their degree-$j$ weight bounds.
The finite-stage filtrations use finite sums and extensions, and
an upper weight bound on cohomology in a fixed degree passes to
extensions.  The resulting bound is uniform in every remaining
index, as asserted.
\end{proof}

\subsection{Completion of the identity}

\begin{theorem}[Plancherel identity]
\label{planch:identity}
For $s\notin\mathbb Q$, the natural transformation \eqref{eq:P6}
is an isomorphism.  Consequently
\[
 J_sR_s\simeq\operatorname{Id}_{D_s},
\]
as asserted in \eqref{eq:P1}.
\end{theorem}

\begin{proof}
On every test chart the family $\{\widehat K_\beta\}$ of
\eqref{planch:tested-cones} is linear in the single frame-ratio Kummer
input and has compatible logarithm lifts.  It therefore satisfies
the finite-stage logarithm and rank hypotheses of
Lemma~\ref{deform:colimit}.  Its cohomology colimit vanishes at
$q=1$ by Lemma~\ref{planch:zero-colimit}, and
Proposition~\ref{planch:weight-estimate} supplies the required
degreewise uniform upper weights.  The colimit lemma and stalk
detection in Lemma~\ref{deform:stalks} therefore give
$\operatorname*{colim}_\beta\widehat K_\beta=0$ at every parameter
that is not a root of unity, in particular at the monodromy of every
irrational $s$.

Fourier transform in the two pre-axes and take $!$-fibres at the
two character parameters $1$.  The pre-axis diagrams are regular
holonomic; Fourier and specialization are computed in the holonomic
Ind-categories on charts.  Away from zero character parameters, the pre-axis
projections do not change them.  We obtain a zero cone after
applying $\mathcal T_s$ to \eqref{eq:P6}.  Its conservativity in
Proposition~\ref{planch:finite-test} proves \eqref{eq:P6} on each
spectator chart.  All arrows used were natural stack morphisms,
so this is the required isomorphism of kernels.  The
pseudo-identity equivalence of Section~\ref{sec:pseudoidentity}
then gives \eqref{eq:P1}.
\end{proof}
